\section{Introducción}

En la contemporaneidad tecnológica, la inmediatez ha dejado de ser una ventaja competitiva para convertirse en un requisito funcional innegociable. Desde la gestión de cámaras de tráfico en las intersecciones más congestionadas de una metrópolis hasta la operación crítica de satélites que deben maniobrar autónomamente para esquivar basura espacial en órbita, el volumen y la velocidad de los datos han superado las capacidades físicas de los sistemas tradicionales. Durante décadas, la industria del software fundamentó sus operaciones en arquitecturas monolíticas y en el patrón de comunicación sincrónico de petición-respuesta (Request-Response). Si bien este modelo fue suficiente para la web estática, hoy se ha convertido en el principal cuello de botella para la innovación en el Internet de las Cosas (IoT). Nos enfrentamos a un escenario donde los sensores y sistemas de vigilancia no pueden darse el "lujo" de esperar segundos a que un servidor centralizado procese su solicitud; la realidad exige una reacción en tiempo real.

Este documento nace de la necesidad imperiosa de explorar cómo la ingeniería de software está respondiendo a este colapso del modelo tradicional. A través de un análisis exhaustivo de veinte investigaciones recientes, evidenciamos un cambio de paradigma radical: la transición desde sistemas acoplados hacia \textit{Arquitecturas Orientadas a Eventos} (EDA) y el uso de \textit{Microservicios Asíncronos}. La premisa fundamental que guía este estudio es que el futuro del software no reside en la acumulación pasiva de información en bases de datos estáticas, sino en la capacidad de capturar, interpretar y reaccionar ante los "eventos" que definen nuestra realidad en el instante exacto en que ocurren.

Sin embargo, la implementación de este nuevo paradigma enfrenta desafíos estructurales profundos. Uno de los problemas más críticos es la ineficiencia de las metodologías actuales de implementación en la nube. Como señalan Shujaat et al. \cite{shujaat2021improvements}, el enfoque convencional de "Cloud-Centric IoT" (CC-IoT) adolece de tiempos de descubrimiento de servicios lentos y una gestión de recursos subóptima. Su investigación destaca que registrar sensores en listas planas sin una estructura lógica (como Árboles de Acceso a Servicios) provoca latencias inaceptables cuando el número de dispositivos escala a millones. Esto sugiere que la solución no es solo "mover datos más rápido", sino organizarlos de manera inteligente antes de que salgan del dispositivo.

A esta complejidad se suma la necesidad de una respuesta contextual. No basta con detectar un evento; el sistema debe saber qué hacer con él según el entorno. Ward et al. \cite{ward2005response} introducen una distinción vital con su concepto de "Arquitectura de Respuesta Dirigida por Eventos" (EDRA). Argumentan que los sistemas actuales son buenos para alertar, pero malos para actuar. Un sistema verdaderamente inteligente debe mantener un "estado de contexto" del usuario o del entorno para decidir si una alerta requiere una acción inmediata o si puede ser ignorada, cerrando la brecha entre la detección técnica y la utilidad del negocio.

Otro obstáculo recurrente que abordaremos es la desconexión semántica y operativa en entornos urbanos. Como bien señalan Phuttharak y Loke \cite{phuttharak2023smartcity}, y reafirmado por Visnjic et al. \cite{visnjic2021decision}, las ciudades inteligentes actuales son, en realidad, archipiélagos de datos desconectados. Los sensores actúan de forma aislada, incapaces de comunicarse entre sí debido a la heterogeneidad de sus formatos y protocolos. Sin una arquitectura de plataforma que estandarice y orqueste estos eventos, la promesa de la *Smart City* se diluye en silos de información inútil. De forma similar, en el ámbito de la salud, veremos cómo la coordinación de ambulancias puede dejar de depender de reportes manuales propensos al error humano, pasando a sistemas de fusión de datos basados en eventos complejos que optimizan tiempos de respuesta vitales \cite{bruns2014fusion}.

La infraestructura subyacente también juega un papel crucial. Park et al. \cite{park2021cloud} demuestran que es posible construir arquitecturas resilientes y de bajo costo (REDA) integrando protocolos ligeros como MQTT con buses de eventos robustos como Kafka, desafiando la noción de que la alta disponibilidad requiere inversiones millonarias en hardware propietario.

No obstante, es preciso aclarar que esta evolución no está exenta de fricciones. La migración de sistemas monolíticos legados a ecosistemas distribuidos es una tarea costosa y técnicamente demandante. La experiencia de Tecgraf/PUC-Rio al modernizar su sistema de monitoreo masivo, documentada por Laigner et al. \cite{laigner2020monolithic}, sirve como una advertencia sobre la complejidad de gestionar la consistencia de datos cuando se abandona la seguridad de las transacciones ACID tradicionales. Además, adoptar ciegamente estas nuevas arquitecturas sin una estrategia clara de observabilidad puede generar "puntos ciegos" peligrosos donde los errores se vuelven imposibles de rastrear, como advierten Lazzari y Farias \cite{lazzari2023hidden}.

A lo largo de este artículo, desglosaremos cómo tecnologías emergentes se integran en este ecosistema. Analizaremos desde el uso de cámaras neuromórficas y algoritmos de visión como Stack-CNN para la prevención de colisiones espaciales \cite{coretti2025neuromorphic}, hasta la detección de infracciones de tráfico mediante YOLOv5 en el borde \cite{perera2024traffic}. El objetivo final es proporcionar una visión holística de cómo la arquitectura de software está evolucionando para sobrevivir en un mundo donde el tiempo de respuesta se mide en milisegundos.